% !TeX program = xelatex
\documentclass[a4paper,10pt]{article}
\usepackage{fontspec}
\usepackage[margin=21mm,headheight=14pt]{geometry}
\usepackage{amsmath,amssymb,booktabs,array}
\usepackage{xcolor,hyperref,fancyhdr,listings}
\setmainfont{Noto Serif CJK JP}
\setsansfont{Noto Sans CJK JP}
\setmonofont{DejaVu Sans Mono}
\XeTeXlinebreaklocale "ja"
\XeTeXlinebreakskip=0pt plus 1pt
\hypersetup{colorlinks=true,urlcolor=blue!50!black,
  pdftitle={第4回課題 解答：Newton法},
  pdfauthor={岩澤大心}}
\pagestyle{fancy}
\fancyhf{}
\lhead{第4回課題：Newton法}
\rhead{25G1017\quad 岩澤大心}
\cfoot{\thepage}
\setlength{\parindent}{1em}
\setlength{\parskip}{0.35em}
\renewcommand{\arraystretch}{1.22}
\lstset{language=Python,basicstyle=\ttfamily\fontsize{8.2}{10.1}\selectfont,
  keywordstyle=\color{blue!55!black},commentstyle=\color{black!55},
  stringstyle=\color{green!35!black},frame=single,rulecolor=\color{black!25},
  backgroundcolor=\color{black!2},columns=fullflexible,keepspaces=true,
  showstringspaces=false,breaklines=true,tabsize=4}
\begin{document}
\begin{center}
  {\LARGE\bfseries 第4回課題 解答}\\[5mm]
  学籍番号：25G1017\qquad 氏名：岩澤大心\\[2mm]
  作成日：2026年10月12日
\end{center}

\section*{課題1：$x^2=2$をNewton法で解く}
$f(x)=x^2-2$，$f'(x)=2x$とおくと，更新式は
\[
 x_{k+1}=x_k-\frac{x_k^2-2}{2x_k},\qquad x_0=2
\]
となる．一般のベクトルの場合は，ヤコビ行列$J$について
$J(\boldsymbol{x}_k)\boldsymbol{\Delta}_k=F(\boldsymbol{x}_k)$を解き，
$\boldsymbol{x}_{k+1}=\boldsymbol{x}_k-\boldsymbol{\Delta}_k$とする．
空欄は\texttt{linalg.solve(Df(x), f(x))}で補い，その後に\texttt{x = x - DFinv\_Fx}を実行する．
逆行列を明示的に作らず，連立一次方程式を直接解く．

\lstinputlisting[firstline=3,lastline=33]{assignment4.py}
\begin{lstlisting}
x = newton([2.0], f1, Df1, cond_fixed)
print(x)  # [1.41421356]
\end{lstlisting}
計算結果は$x\simeq1.414213562373095$であり，正の解$\sqrt{2}$に一致する．
更新後に終了判定を行うため，$i=0,\ldots,20$の\texttt{i == 20}は
計21回の更新に相当する．残差$|x^2-2|$は約$4.44\times10^{-16}$であった．
配布コードの1要素配列は\texttt{x[0]}として参照する．
\texttt{history}は履歴保存用，\texttt{max\_iter}は未収束時の無限ループ防止用である．

\newpage
\section*{課題2：更新量による終了条件}
更新量の全成分の絶対値を取り，最大値が$10^{-14}$以下になったら終了する．
すなわち，
\[
  \|\boldsymbol{\Delta}_k\|_\infty
  =\max_j|(\boldsymbol{\Delta}_k)_j|\le 10^{-14}
\]
を判定する．課題1の回数による判定を，次の\texttt{cond}に置き換える．
\lstinputlisting[firstline=36,lastline=37]{assignment4.py}
\begin{lstlisting}
history = []
x = newton([2.0], f1, Df1, cond, history=history)
print(x)
print(len(history))
\end{lstlisting}
出力は\texttt{[1.41421356]}，更新回数は6回となった．
更新後の値と，その更新に用いた修正量の絶対値を表に示す．

\begin{center}
\begin{tabular}{crr}
\toprule
更新回数 & 更新後の$x$ & $|\Delta|$\\
\midrule
1 & 1.500000000000000 & $5.000000000\times10^{-1}$\\
2 & 1.416666666666667 & $8.333333333\times10^{-2}$\\
3 & 1.414215686274510 & $2.450980392\times10^{-3}$\\
4 & 1.414213562374690 & $2.123899820\times10^{-6}$\\
5 & 1.414213562373095 & $1.594742910\times10^{-12}$\\
6 & 1.414213562373095 & $1.570092459\times10^{-16}$\\
\bottomrule
\end{tabular}
\end{center}
5回目の更新量は$10^{-14}$より大きく，6回目で条件を満たす．
課題1と同じ精度の解を，固定回数より少ない更新で得られた．

\subsection*{結果の確認}
\[
  x\simeq1.414213562373095,\qquad
  |x^2-2|\simeq4.44\times10^{-16}.
\]
終了条件は\texttt{f(x)}の値ではなく，指定された
\texttt{DFinv\_Fx}の最大絶対値に対する条件である．
残差は解の確認用に別途計算した．例えば，修正量が
$(5\times10^{-15},-2\times10^{-14})$なら最大絶対値は
$2\times10^{-14}$であり，終了しない．

\newpage
\section*{課題3：非線形連立方程式}
与えられた方程式を
\[
 F(x,y)=
 \begin{pmatrix}
  2x^2-e^y\\
  -3xy^2-y^3-3
 \end{pmatrix}
 =\boldsymbol{0}
\]
とおく．各成分を$x,y$で偏微分したヤコビ行列は
\[
 J(x,y)=
 \begin{pmatrix}
  4x & -e^y\\
  -3y^2 & -6xy-3y^2
 \end{pmatrix}
\]
である．課題1の\texttt{newton}と，課題2の\texttt{cond}をそのまま利用する．

\lstinputlisting[firstline=40,lastline=49]{assignment4.py}
\begin{lstlisting}
x = np.zeros(2)
x[0] = 4.0
x[1] = 5.0
history = []
x = newton(x, f3, Df3, cond, history=history)
print(x)
print(len(history))
print(np.max(np.abs(f3(x))))
\end{lstlisting}

\subsection*{計算結果}
指定された初期値$(x_0,y_0)=(4,5)$から，40回の更新で
\[
 \boxed{x\simeq-1.214687810300159,\qquad
        y\simeq1.082121375602253}
\]
を得た．小数点以下8桁に丸めると
$(x,y)=(-1.21468781,\ 1.08212138)$となり，課題文の近似解と一致する．

最終更新量の最大絶対値は約$1.18\times10^{-16}$，残差は
\[
 \|F(x,y)\|_\infty\simeq4.44\times10^{-16}
\]
であった．以上より，終了条件を満たし，元の2式も十分な精度で成立している．
この初期値では20回固定の更新では収束しないため，
課題2の更新量による終了条件を用いた．

\newpage
\section*{課題4：初期値を変えて別の解を探す}
課題3と同じ関数・ヤコビ行列・終了条件を使い，初期値を変更した．
課題3以外に得られた近似解は，次の2種類である．
\begin{align*}
 \text{解A：}\quad &(x,y)\simeq
 (0.288166073001912,\ -1.795289464785586),\\
 \text{解B：}\quad &(x,y)\simeq
 (-0.401282369883953,\ -1.133032689807632).
\end{align*}
これらは課題3の$y>0$の解とは異なり，いずれも$y<0$である．

\begin{lstlisting}
for initial in ([1.0, -2.0], [4.0, -5.0], [1.0, 1.0],
                [-1.0, -1.0], [-4.0, -5.0], [2.0, -2.0]):
    history = []
    root = newton(initial, f3, Df3, cond, history=history)
    print(initial, root, len(history),
          np.max(np.abs(f3(root))))
\end{lstlisting}

\begin{center}
\begin{tabular}{cccr}
\toprule
初期値$(x_0,y_0)$ & 収束先 & 更新回数 & $\|F(x,y)\|_\infty$\\
\midrule
$(1,-2)$  & 解A & 10 & $4.44\times10^{-16}$\\
$(4,-5)$  & 解A &  9 & $4.44\times10^{-16}$\\
$(1,1)$   & 解A & 12 & $4.44\times10^{-16}$\\
$(-1,-1)$ & 解B &  7 & $4.44\times10^{-16}$\\
$(-4,-5)$ & 解B & 10 & $8.88\times10^{-16}$\\
$(2,-2)$  & 解B & 12 & $8.88\times10^{-16}$\\
\bottomrule
\end{tabular}
\end{center}

\subsection*{確認と考察}
課題3の解と合わせて，今回の計算では3種類の解を確認した．
例えば$(1,-2)$と$(2,-2)$は，初期値の符号が同じでも異なる解へ収束した．
Newton法の収束先は初期値に依存し，初期値の符号だけでは決まらない．
また，特異なヤコビ行列や発散が生じる初期値もあるため，
終了条件の成立に加えて残差も確認した．

第2式から$y=0$は不可能であり，$y\ne0$では
\[
  x=-\frac{y}{3}-\frac{1}{y^2}
\]
と整理できる．得られた3種類の解は，この関係と第1式
$2x^2=e^y$の両方を数値的に満たす．

\subsection*{再実行方法}
同梱の\texttt{assignment4.py}を\texttt{python assignment4.py}で実行すると，
課題1--4の初期値・近似解・更新回数・最終更新量・残差を表示する．
実行環境はPython 3，NumPy 2.3.5，SciPy 1.17.0である．
表の残差は丸め表示前の計算値から求めた．

\end{document}
